% !TEX TS-program = lualatex+se
% !TEX encoding = UTF-8

\documentclass[oneside, 12pt, a4paper]{book}
%\usepackage[letterpaper,showframe]{geometry}
\usepackage{fullpage} % redukcja marginesów
\usepackage{geometry}
\geometry{
	left=12mm,
	right=12mm,
	top=12mm,
	bottom=12mm,
	marginparwidth=0mm,
	marginparsep=0mm
	}
\usepackage{gregoriotex}
\gresetlinecolor{gregoriocolor}
\gresetheadercapture{annotation}{greannotation}{string}
% Format annotation above initial
\grechangestyle{annotation}{\small\bfseries}%
\grechangedim{intersyllablespacenotes}{0.18cm}{scalable}%

\usepackage{fullpage}
\usepackage{multicol}
\usepackage{setspace} % odstępy między liniami

%  Konfiguracja kolumn
\usepackage{paracol}
\setcolumnwidth{14cm,4cm}
\setlength{\columnsep}{0.4cm}
% \setlength{\columnseprule}{0.4pt}
% \colseprulecolor{red}

\usepackage{tcolorbox}
\newtcolorbox{nutybox}{size=minimal, colback=white, colframe=red, arc=0mm, width=\linewidth, halign=left, valign=center}
\newtcolorbox{tlumbox}{size=minimal, colback=white, colframe=red, arc=0mm, width=\linewidth, halign=left, valign=top}

\usepackage{Tabbing}
\usepackage{longtable}
\usepackage[savepos]{zref}

\let\grelocalleftbox\localleftbox
\let\grelocalrightbox\localrightbox

\usepackage[latin,polish]{babel}

\usepackage{fontspec}
\defaultfontfeatures{Ligatures=TeX}
\setmainfont{Libertinus Serif}

\usepackage{libertine} % ust. czcionki
\pagestyle{empty}	% usuwanie nagłówków/stopek z dokumentu
\usepackage[indent=0pt]{parskip} % ustawianie paragrafów
%\setstretch{1.5} 

\usepackage{graphicx} % umożliwia import obrazów
\graphicspath{{img/}} % katalog obrazów
\usepackage{pgfornament}% ornamenty graficzne

\usepackage{titlesec} % ustawianie nagłówków
\titlespacing*{\chapter}{0pt}{0pt}{0pt}	
\titleformat{\chapter}[hang]{\filcenter\LARGE\scshape\bfseries}{}{0pt}{} % format nagłówków
\titlespacing*{\section}{0pt}{0pt}{0mm}	
\titleformat{\section}[hang]{\filcenter\large\scshape\bfseries}{}{0pt}{} % format nagłówków

%\newfontface\vectis{VectisMonumentalW01-Regular}
%\def\greinitialformat#1{\fontsize{20}{20}\vectis #1}

% --- komendy własne ---
\newcommand{\red}[1]{\textcolor{red}{#1}}
\newcommand{\black}[1]{\textcolor{black}{#1}}

\newcommand{\ornament}[1]{
	\begin{center}
		\vspace{25mm}
		\pgfornament[width=0.6cm]{#1}
		\newpage
	\end{center}}

\begin{document}

\setlength{\parindent}{0pt}
\gresetbracerendering{font}

\pagenumbering{gobble} % nie jest wyprowadzany żaden numer strony, chociaż numer jest nadal resetowany. Odniesienia do tej strony również są puste.

\begin{spacing}{0.7}
\begin{center}
	\section*{A solis ortus cardine}
	\small\textit{Liber antiphonarius, 1960, p. 266}
	\vspace{5mm}
\end{center}	
\end{spacing}
\begin{paracol}{2}
	\ensurevspace{3.5cm}
	\begin{nutybox}\greannotation{Hymn.} \greannotation{3.} \gabcsnippet{(c4) A(d) so(e')lis(f) or(gh)tus(d) cár(eg)di(fe)ne(e.) (;) Ad(g) u(hj)sque(j) ter(ji)rae(hg) lí(hi)mi(i)tem,(i.) (:) Chri(h)stum(hj!kv) ca(j)ná(ji)mus(hg) Prín(h)ci(gf)pem,(e_d/f!gh.) (;) Na(d)tum(e') Ma(f)rí(gh)a(gf) Vír(g')gi(f)ne.(e.) (::)}
	\end{nutybox}
	\switchcolumn
	\begin{tlumbox}
	\vspace{3.5mm} 
	Od kresu wschodu słońca aż do krańca ziemi opiewajmy Chrystusa Króla, narodzonego z~Marji Dziewicy.
	\end{tlumbox}
  	\ensurevspace{3.5cm}
	\switchcolumn*
	\begin{nutybox}
	\gresetinitiallines{0}\gabcsnippet{(c4) <c>2.</c> Be(d)á(e')tus(f) au(gh)ctor(d) saé(eg)cu(fe)li(e.) (;) Ser(g)ví(hj)le(j) cor(ji)pus(hg) ín(hi)du(i)it:(i.) (:) Ut(h) car(hj!kv)ne(j) car(ji)nem(hg) lí(h)be(gf)rans,(e_d/f!gh.) (;) Ne(d) pér(e')de(f)ret(gh) quos(gf) cón(g')di(f)dit.(e.) (::)}
	\end{nutybox}
	\switchcolumn
	\begin{tlumbox}
	\vspace{3.5mm} 
	Błogosławiony Stwórca świata przybrał ciało sługi, by ciało ciałem wyzwalając, nie stracił (tych), których stworzył.
	\end{tlumbox}
	\ensurevspace{3.5cm}
	\switchcolumn*
	\begin{nutybox}
	\gresetinitiallines{0}\gabcsnippet{ (c4) <c>3.</c> Ca(d)stae(e') Pa(f)rén(gh)tis(d) ví(eg)sce(fe)ra(e.) (;) Cae(g)lé(hj)stis(j) in(ji)trat(hg) grá(hi)ti(i)a:(i.) (:) Ven(h)ter(hj!kv) pu(j)él(ji)lae(hg) bá(h)ju(gf)lat(e_d/f!gh.) (;) Se(d)cré(e')ta,(f) quae(gh) non(gf) nó(g')ve(f)rat.(e.) (::)}
	\end{nutybox}
	\switchcolumn
	\begin{tlumbox}
	\vspace{3.5mm} 
	Czystej Rodzicielki wnętrzności nawiedza łaska niebieska: żywot Dziewicy nosi tajemnice, których nie poznała.
	\end{tlumbox}
	\ensurevspace{3.5cm}
	\switchcolumn*
	\begin{nutybox}
	\gresetinitiallines{0}\gabcsnippet{ (c4) <c>4.</c> Do(d)mus(e') pu(f)dí(gh)ci(d) pé(eg)cto(fe)ris(e.) (;) Tem(g)plum(hj) re(j)pén(ji)te(hg) fit(hi) De(i)i:(i.) (:) In(h)tá(hj!kv)cta(j) né(ji)sci(hg)ens(h) vi(gf)rum,(e_d/f!gh.) (;) Con(d)cé(e')pit(f) al(gh)vo(gf) Fí(g')li(f)um.(e.) (::)}
	\end{nutybox}
	\switchcolumn
	\begin{tlumbox}
	\vspace{3.5mm} 
	Przybytek wstydliwych piersi niespodzianie staje się świątynią Boga: nienaruszona, nie znając męża, poczęła Syna w~żywocie.
	\end{tlumbox}
	\ensurevspace{3.5cm}
	\switchcolumn*
	\begin{nutybox}
	\gresetinitiallines{0}\gabcsnippet{(c4) <c>5.</c> E(d)ní(e')ti(f)tur(gh) pu(d)ér(eg)pe(fe)ra(e.) (;) Quem(g) Gá(hj)bri(j)el(ji) prae(hg)dí(hi)xe(i)rat,(i.) (:) Quem(h) ven(hj!kv)tre(j) ma(ji)tris(hg) gé(h)sti(gf)ens(e_d/f!gh.) (;) Bap(d)tí(e')sta(f) clau(gh)sum(gf) sén(g')se(f)rat.(e.) (::)}
	\end{nutybox}
	\switchcolumn
	\begin{tlumbox}
	\vspace{3.5mm} 
	Wydaje rodzicielka (na świat Tego), którego Gabrjel przepowiedział, którego w~żywocie Matki zamkniętego, radując się, przeczuł Chrzciciel.
	\end{tlumbox}
	\ensurevspace{3.5cm}
	\switchcolumn*
	\begin{nutybox}
	\gresetinitiallines{0}\gabcsnippet{(c4) <c>6.</c> Foe(d)no(e') ja(f)cé(gh)re(d) pér(eg)tu(fe)lit,(e.) (;) Prae(g)sé(hj)pe(j) non(ji) a(hg)bhór(hi)ru(i)it:(i.) (:) Et(h) la(hj!kv)cte(j) <i>mó</i>()di(ji)co(hg) pa(h)stus(gf) est,(e_d/f!gh.) (;) Per(d) quem(e') nec(f) a(gh)les(gf) é(g')su(f)rit.(e.) (::)}
	\end{nutybox}
	\switchcolumn
	\begin{tlumbox}
	\vspace{3.5mm} 
	Znosił (cierpliwie) leżenie na sianie, żłóbkiem nie pogardził i~odrobiną mleka karmił się (Ten), dzięki któremu nawet ptak nie łaknie.
	\end{tlumbox}
	\ensurevspace{3.5cm}
	\switchcolumn*
	\begin{nutybox}
	\gresetinitiallines{0}\gabcsnippet{(c4) <c>7.</c> Gau(d)det(e') cho(f)rus(gh) cae(d)lé(eg)sti(fe)um,(e.) (;) Et(g) An(hj)ge(j)li(ji) ca(hg)nunt(hi) De(i)o:(i.) (:) Pa(h)lám(hj!kv)que(j) fit(ji) pa(hg)stó(h)ri(gf)bus(e_d/f!gh.) (;) Pa(d)stor,(e') Cre(f)á(gh)tor(gf) ó(g')mni(f)um.(e.) (::)}
	\end{nutybox}
	\switchcolumn
	\begin{tlumbox}
	\vspace{3.5mm} 
	Raduje się chór niebian, i~Aniołowię śpiewają Bogu; jawnym staje się pasterzom Pasterz, Stwórca wszystkich.
	\end{tlumbox}
	\ensurevspace{3.5cm}
	\switchcolumn*
	\begin{nutybox}
	\gresetinitiallines{0}\gabcsnippet{(c4) <c>8.</c> Je(d)su,(e') ti(f)bi(gh) sit(d) gló(eg)ri(fe)a,(e.) (;) Qui(g) na(hj)tus(j) es(ji) de(hg) Vír(hi)gi(i)ne,(i.) (:) Cum(h) Pa(hj!kv)tr<i>e</i>() et(j) al(ji)mo(hg) Spí(h)ri(gf)tu,(e_d/f!gh.) (;) In(d) sem(e')pi(f)tér(gh)na(gf) saé(g')cu(f)la.(e.) (::) A(efe)men.(de..) (::)}
	\end{nutybox}
	\switchcolumn
	\begin{tlumbox}
	\vspace{3.5mm} 
	Jezu, Tobie niechaj będzie chwała, który narodziłeś się z Dziewicy, oraz Ojcu i~dobroczynnemu Duchowi, po nieskończone wieki. Amen.\footnotemark
	\end{tlumbox}
\end{paracol}	
\footnotetext{tłum. Ks. Bronisław Gładysz, "Hymny brewiarza rzymskiego oraz patronału polskiego. Przekład i objaśnienia", Poznań 1933}
\end{document}